\begin{center}
\textbf{\large{Tóm tắt}}
\end{center}

\addcontentsline{toc}{chapter}{Tóm tắt}

\begin{small}

Kiểm soát truy cập giới hạn những thao tác mà mỗi người dùng được phép thực hiện trên REST API. Các kỹ thuật kiểm thử tự động hiện nay chủ yếu tập trung vào lỗi phân quyền phi trạng thái, như lỗi ở mức đối tượng (BOLA) hoặc mức chức năng (BFLA). Lớp lỗ hổng phân quyền có trạng thái (Stateful Broken Access Control) ít được các kỹ thuật này kiểm tra trực tiếp. Lỗi này xuất hiện khi quyền hạn của người dùng bị thu hồi trong vận hành, nhưng hệ thống vẫn chấp nhận các phiên làm việc cũ.

Khóa luận đề xuất giải pháp tự động phát hiện lớp lỗ hổng này bằng kiểm thử biến hình trên công cụ EvoMaster. Giải pháp dựa trên quan hệ MR-A1: nếu một thao tác cần đúng quyền vừa bị thu hồi và đã thành công trước khi tài khoản bị hạ quyền, thì khi gửi lại bằng cùng token sau đó, thao tác phải bị từ chối. Mỗi kịch bản kiểm thử gồm bộ ba thao tác: thao tác gốc, thao tác hạ quyền và thao tác kiểm chứng. Mô hình Naive Bayes tự động nhận diện thao tác hạ quyền từ đặc tả OpenAPI, nên người dùng không phải tự khai báo thao tác này.

Thực nghiệm gồm 98 lượt chạy ở chế độ hộp đen trên bốn ứng dụng mục tiêu. reports-api và vampi-seeded mang lỗi tự gieo; pyLoad mang lỗ hổng \mbox{CVE-2026-41133}; artifact-keeper mang một lỗi mà nhà phát triển đã vá. Trên các biến thể có lỗ hổng, công cụ phát hiện vi phạm ở 19/20 lượt với reports-api, 6/20 lượt với vampi-seeded, 5/5 lượt với pyLoad và 4/5 lượt với artifact-keeper. Trong 48 lượt chạy trên các biến thể đã vá, công cụ dựng được bộ ba ở 30 lượt và không báo động giả ở lượt nào.

\vspace*{1cm}
{\raggedright\textbf{Từ khóa}: REST API, kiểm soát truy cập, phân quyền có trạng thái, kiểm thử biến hình, EvoMaster\par}

\end{small}
